\documentclass[11pt]{amsart}
\usepackage{amsmath,amssymb,amsthm,booktabs}
\usepackage[margin=1.15in]{geometry}
\usepackage[hidelinks]{hyperref}
\newtheorem{theorem}{Theorem}[section]
\newtheorem{lemma}[theorem]{Lemma}
\newtheorem{proposition}[theorem]{Proposition}
\newtheorem{corollary}[theorem]{Corollary}
\newtheorem{conjecture}[theorem]{Conjecture}
\theoremstyle{remark}
\newtheorem{remark}[theorem]{Remark}
\newcommand{\PP}{\mathbb{P}}
\newcommand{\E}{\mathbb{E}}
\newcommand{\vE}{\vec{E}}
\newcommand{\doi}[1]{\href{https://doi.org/#1}{doi:\nolinkurl{#1}}}
\setlength{\emergencystretch}{2.5em}

\title[Non-backtracking walks on graphs of bounded degree ratio]{A birthday bound for non-backtracking walks on graphs
of bounded degree ratio}
\author{Yitzchak Shmalo}
\address{Einstein Institute of Mathematics, The Hebrew University of Jerusalem, Jerusalem, Israel}
\email{yspennstate@gmail.com}
\date{November 2026}
\subjclass[2020]{05C81, 60J10, 60C05}
\keywords{non-backtracking random walk, self-intersection, birthday paradox, irregular graphs}

\raggedbottom
\begin{document}
\begin{abstract}
Let $G$ be a finite simple graph on $n$ vertices with minimum degree $m\ge3$ and maximum degree $M$, and put
$Q=M/(m-1)$. A non-backtracking random walk of length $k$ on $G$, started along any directed edge, is
vertex-simple with probability at most $256\,Q(1+Q)\,n/k^2$. Hence the conjecture of Alon and Peres, that a
non-backtracking walk of length $C(\varepsilon)\sqrt n$ self-intersects with probability at least $1-\varepsilon$,
holds with $C(\varepsilon)$ depending only on $\varepsilon$ and a bound for $M/m$; for regular graphs of degree at
least three one may take $C(\varepsilon)=31\varepsilon^{-1/2}$ for every degree. The expected self-intersection time
is at most $32\sqrt{Q(1+Q)n}+2$, which removes the logarithmic factor from the bound of Bencs et al.\ for graphs of
minimum degree at least three and bounded maximum degree. The proof adapts Dozier's argument for regular graphs:
departures are weighted by the inverse branching at their vertex, and the deviating local derangement is resampled.
\end{abstract}
\maketitle

\section{Introduction}\label{sec:intro}

Let $G=(V,E)$ be a finite simple graph with $n=|V|$ vertices, write $N(v)$ for the neighbourhood and $d(v)=|N(v)|$
for the degree of a vertex $v$, and assume that the minimum degree $m$ is at least $3$. Let $M$ be the maximum degree
and $\vE=\{(u,v):\{u,v\}\in E\}$ the set of directed edges. The non-backtracking walk started along
$\delta=(x_0,x_1)\in\vE$ is the sequence $X_0=x_0$, $X_1=x_1$, and, for $t\ge1$, $X_{t+1}$ chosen uniformly from
$N(X_t)\setminus\{X_{t-1}\}$, independently of the past given $(X_{t-1},X_t)$. For $k\ge1$ put
\[
p_k(\delta)=\PP\bigl(X_0,X_1,\dots,X_k\text{ are distinct}\bigr).
\]
A walk of length $k$ has $k+1$ vertices, so $p_k(\delta)=0$ for $k\ge n$.

Dozier~\cite{Doz26} proved that for fixed $d\ge3$ and $\varepsilon>0$ there is $C=C(d,\varepsilon)$ such that
$p_k(\delta)<\varepsilon$ for every $d$-regular graph on $n$ vertices, every $\delta$ and every $k>C\sqrt n$. He stated
the following conjecture of Alon and Peres, of which this is the fixed-degree case.

\begin{conjecture}[Alon--Peres; {\cite[Conjecture~1.5]{Doz26}}]\label{conj:AP}
For every $\varepsilon>0$ there is $C=C(\varepsilon)$ such that if $G$ is a graph on $n$ vertices with all degrees at
least $3$ and $k>C\sqrt n$, then $p_k(\delta)<\varepsilon$ for every directed edge $\delta$.
\end{conjecture}

Closely related, Bencs, Goldberg, Jenssen, Jerrum, Pete, Regts and Yin~\cite{BGJ26} studied the self-intersection
time of the walk started at a vertex $u$, whose first step is uniform on $N(u)$,
\[
T(G,u)=\min\{t\ge2: X_t\in\{X_0,\dots,X_{t-1}\}\}.
\]
They proved $\E\,T(G,u)=O_\Delta(\sqrt n\log n)$ for connected graphs with minimum degree at least $3$ and maximum
degree $\Delta$~\cite[Theorem~1]{BGJ26}, improved this to $O(\sqrt n)$ for regular graphs with a uniform spectral
gap~\cite[Theorem~2]{BGJ26}, showed that $\sqrt n$ is the right order on certain regular expanders~\cite[Theorem~3]{BGJ26},
and asked whether $\E\,T(G,u)=O(\sqrt n)$ for every graph with minimum degree at least three~\cite[Question~4]{BGJ26}.
For multigraphs the answer is negative; an example follows that question in~\cite{BGJ26}. The non-backtracking walk
is a Markov chain on directed edges, and its mixing on regular graphs was studied by Alon, Benjamini, Lubetzky and
Sodin~\cite{ABLS07} and by Lubetzky and Peres~\cite{LP16}; Theorem~2 of~\cite{BGJ26} combines the latter with the
birthday paradox for Markov chains of Kim, Montenegro, Peres and Tetali~\cite[Theorem~3.6]{KMPT10}.

Put
\[
Q=\frac{M}{m-1}.
\]

\begin{theorem}\label{thm:main}
Let $G$ be a finite simple graph on $n$ vertices with minimum degree $m\ge3$ and maximum degree $M$. For every
directed edge $\delta$ and every $k\ge1$,
\[
p_k(\delta)\le\min\Bigl\{1,\ \frac{256\,Q(1+Q)\,n}{k^2}\Bigr\}.
\]
\end{theorem}

The constant $256$ is not optimised; the numerical checks described at the end of the paper suggest that
$\max_{\delta,k}p_k(\delta)k^2/n$ is of order one on the graphs tested.

\begin{corollary}\label{cor:constants}
Let $0<\varepsilon<1$ and $D\ge1$. If $M/m\le D$ and $k>32D\sqrt{n/\varepsilon}$, then $p_k(\delta)<\varepsilon$ for
every $\delta$. If $M\le\Delta$ and $k>8\sqrt{\Delta(\Delta+2)n/\varepsilon}$, then $p_k(\delta)<\varepsilon$ for every
$\delta$. In particular, if $G$ is regular of any degree $d\ge3$ and $k>31\sqrt{n/\varepsilon}$, then
$p_k(\delta)<\varepsilon$.
\end{corollary}

Thus Conjecture~\ref{conj:AP} holds on every class of graphs with bounded ratio of maximum to minimum degree, with a
constant depending only on $\varepsilon$ and the bound for the ratio. For regular graphs this gives Dozier's theorem
with a constant that does not depend on the degree.

For $\delta\in\vE$ let $\tau(\delta)=\min\{t\ge2:X_t\in\{X_0,\dots,X_{t-1}\}\}$ for the walk started along $\delta$.

\begin{corollary}\label{cor:expectation}
For every directed edge $\delta$, $\E\,\tau(\delta)\le32\sqrt{Q(1+Q)n}+2$. Consequently, for every vertex $u$,
$\E\,T(G,u)\le32\sqrt{Q(1+Q)n}+2$; if $M\le\Delta$ this is at most $16\sqrt{\Delta(\Delta+2)n}+2$, and if
$M/m\le D$ it is at most $62D\sqrt n+2$.
\end{corollary}

For minimum degree at least three and maximum degree at most $\Delta$ this improves~\cite[Theorem~1]{BGJ26} from
$O_\Delta(\sqrt n\log n)$ to $O(\Delta\sqrt n)$, without assuming that $G$ is connected, and by~\cite[Theorem~3]{BGJ26}
the order $\sqrt n$ cannot be improved. It answers~\cite[Question~4]{BGJ26} for graphs of bounded degree ratio.

The proof follows the argument of~\cite{Doz26}. A random plan fixes, at every vertex, how the walk leaves it as a
function of how it arrived, so that the walk becomes a deterministic function of the plan; the measure on plans is
chosen so that the plan walk is vertex-simple with probability exactly $p_k(\delta)$ (Section~\ref{sec:plans}). From a
plan whose walk is simple, one deviates at an interior vertex of the walk. Deviations that never return to the walk are
few because their continuations are disjoint, and deviations that return can be paired by reversing the excursion, so
that a fixed proportion of all deviations self-intersect (Section~\ref{sec:departures}). Since an excursion is short,
each self-intersecting plan arises from few simple ones (Section~\ref{sec:flow}). Two changes make this work when the
degrees vary. Each deviation at a vertex $v$ is given weight $1/(d(v)-1)$, the probability that the walk takes a given
exit at $v$; excursion reversal then changes weights by a factor at most $Q$. And the local derangement at the
deviating vertex is resampled uniformly among those that take the chosen exit, which makes the transport of mass from
simple plans to self-intersecting plans exact and avoids a factor equal to the number of derangements of $M$ objects.
The degree ratio enters at three places: the number of non-returning deviations, the number of deviations returning to
the last vertex of the walk, and the change of weight under reversal (Remark~\ref{rem:ratio}).

\section{Plans}\label{sec:plans}

For $d\ge2$ let $R_d$ be the number of derangements (fixed-point-free permutations) of a set with $d$ elements. A
\emph{plan} is a family $\rho=(\pi_v)_{v\in V}$ in which $\pi_v$ is a derangement of $N(v)$. A plan defines a map
$\phi_\rho:\vE\to\vE$, $\phi_\rho(u,v)=(v,\pi_v(u))$. At each vertex $v$ it maps the $d(v)$ directed edges entering $v$
bijectively onto the $d(v)$ directed edges leaving $v$, so $\phi_\rho$ is a permutation of $\vE$, and since $\pi_v$ has
no fixed point, $\phi_\rho(u,v)\ne(v,u)$. The \emph{walk} of $\rho$ from $\delta$ is the sequence of directed edges
$\phi_\rho^{\,t}(\delta)$, $t\ge0$, and its vertices $x_0,x_1,\dots$, where $\phi_\rho^{\,t}(\delta)=(x_t,x_{t+1})$.

Let $\mu$ be the uniform probability measure on plans, that is, the product of the uniform measures on the derangements
of each $N(v)$. Every plan has $\mu$-mass $\prod_{v\in V}R_{d(v)}^{-1}$. Fix $\delta$ and $k$, and let $S$ be the set of
plans whose walk from $\delta$ has $k+1$ distinct vertices $x_0,\dots,x_k$.

\begin{lemma}\label{lem:derangement}
Let $|N|=d\ge2$ and let $a\ne b$ be elements of $N$. Exactly $R_d/(d-1)$ derangements $\sigma$ of $N$ satisfy
$\sigma(a)=b$.
\end{lemma}

\begin{proof}
A permutation $\eta$ of $N$ fixing $a$ maps derangements to derangements by $\sigma\mapsto\eta\sigma\eta^{-1}$, and
$\eta\sigma\eta^{-1}(a)=\eta(\sigma(a))$. The permutations fixing $a$ act transitively on $N\setminus\{a\}$, so the
number of derangements with $\sigma(a)=b$ is the same for all $d-1$ choices of $b$, and these numbers add up to $R_d$.
\end{proof}

\begin{lemma}\label{lem:simple}
$\mu(S)=p_k(\delta)$.
\end{lemma}

\begin{proof}
Let $y_0,\dots,y_k$ be distinct vertices forming a path with $(y_0,y_1)=\delta$. The walk of $\rho$ from $\delta$
follows this path exactly when $\pi_{y_i}(y_{i-1})=y_{i+1}$ for $1\le i\le k-1$. These conditions concern distinct
vertices, hence are independent under $\mu$, and by Lemma~\ref{lem:derangement} the $i$-th has probability
$1/(d(y_i)-1)$. The probability that the non-backtracking walk follows the path is the same product
$\prod_{i=1}^{k-1}1/(d(y_i)-1)$. Summing over all such paths gives the lemma.
\end{proof}

The walk of a plan and the non-backtracking walk have the same law up to the first repeated vertex, but not after it;
only the identity of Lemma~\ref{lem:simple} is used.

\section{Departures}\label{sec:departures}

Fix $\delta$, $k$ and an integer $L\ge1$. Let $\rho\in S$, let $P=(x_0,\dots,x_k)$ be its walk and $V(P)$ the set of
its vertices. A \emph{departure} of $\rho$ is a pair $(i,\beta)$ with $1\le i\le k-1$ and $\beta=(x_i,z)\in\vE$ with
$z\notin\{x_{i-1},x_{i+1}\}$. Its \emph{weight} is
\[
w_\rho(i)=\frac1{d(x_i)-1}.
\]
At index $i$ there are $d(x_i)-2$ departures, so the total weight of the departures of $\rho$ is
\begin{equation}\label{eq:total}
\sum_{i=1}^{k-1}\frac{d(x_i)-2}{d(x_i)-1}\ge\frac{k-1}2.
\end{equation}

The \emph{excursion} of the departure $(i,\beta)$ is the sequence of directed edges $f_1=\beta$,
$f_{s+1}=\phi_\rho(f_s)$, stopped at the first $s$ for which the head of $f_s$ lies in $V(P)$, or at $s=L$ if there is
no such $s\le L$. Write $f_s=(y_{s-1},y_s)$ for $1\le s\le r$, where $r\le L$ is the length of the excursion, so
$y_0=x_i$ and $y_1,\dots,y_{r-1}\notin V(P)$. The departure is \emph{returning} if $y_r\in V(P)$; then $y_r=x_l$ for a
unique $l$, its \emph{return index}. Otherwise it is \emph{non-returning}, $r=L$ and $y_1,\dots,y_L\notin V(P)$. The
departure is \emph{certified} if the sequence $x_0,\dots,x_i,y_1,\dots,y_r$ has a repeated vertex, that is, if
$y_1,\dots,y_r$ are not distinct or $y_r=x_l$ with $l\le i$.

\begin{lemma}\label{lem:certified}
Let $(i,\beta)$, $\beta=(x_i,z)$, be a certified departure of $\rho\in S$, and let $\psi$ be a plan that agrees with
$\rho$ at every vertex other than $x_i$ and whose derangement at $x_i$ maps $x_{i-1}$ to $z$. Then the walk of $\psi$
from $\delta$ has a repeated vertex among its first $i+L+1$ vertices. In particular $\psi\notin S$ if $i\le k-L$.
\end{lemma}

\begin{proof}
The walk of $\psi$ visits $x_0,\dots,x_i$, because the derangements at $x_1,\dots,x_{i-1}$ are those of $\rho$; its next
directed edge is $\beta$. The vertices $y_1,\dots,y_{r-1}$ lie outside $V(P)$, so they differ from $x_i$ and the
derangements of $\psi$ there are those of $\rho$. Hence the walk of $\psi$ continues along $y_1,\dots,y_r$, and its
first $i+r+1$ vertices are $x_0,\dots,x_i,y_1,\dots,y_r$, which contain a repeated vertex. Since $r\le L$, the claim
follows.
\end{proof}

\begin{lemma}\label{lem:nonreturning}
The departures of $\rho\in S$ that are non-returning and not certified have total weight at most $Qn/L$.
\end{lemma}

\begin{proof}
Let $(i,\beta)$ be such a departure, with excursion $f_1,\dots,f_L$. These directed edges are distinct: if $f_s=f_t$
with $s<t$, their tails would coincide, but the tail of $f_t$ is $y_{t-1}$, which lies outside $V(P)$ and so differs
from $x_i$ when $s=1$, and which differs from $y_{s-1}$ when $s\ge2$ because the departure is not certified.

The excursions of two different such departures $(i,\beta)$ and $(i',\beta')$, with edges $f_s$ and $f'_s$, are
disjoint. Suppose that $f_s=f'_t$ with $s\le t$. Since $f_{s'+1}=\phi_\rho(f_{s'})$, $f'_{t'+1}=\phi_\rho(f'_{t'})$ and
$\phi_\rho$ is injective, $f_{s-a}=f'_{t-a}$ for $0\le a<s$, and in particular $\beta=f_1=f'_{t-s+1}$. If $t>s$, the tail
of $f'_{t-s+1}$ is a vertex $y'_{t-s}$ outside $V(P)$, while the tail of $\beta$ is $x_i\in V(P)$. So $t=s$ and
$\beta=\beta'$, and since $P$ is simple the tail $x_i$ of $\beta$ determines $i$; the two departures coincide.

There are $\sum_v d(v)\le nM$ directed edges, so there are at most $nM/L$ such departures, each of weight at most
$1/(m-1)$.
\end{proof}

\begin{lemma}\label{lem:terminal}
The returning departures of $\rho\in S$ with return index $k$ have total weight at most $Q$.
\end{lemma}

\begin{proof}
Assign to such a departure the last directed edge $f_r$ of its excursion, which enters $x_k$. If two such departures
have $f_r=f'_{r'}$ with $r\le r'$, the argument of Lemma~\ref{lem:nonreturning} gives $\beta=f'_{r'-r+1}$, then
$r=r'$ and $\beta=\beta'$. So there are at most $d(x_k)\le M$ such departures, each of weight at most $1/(m-1)$.
\end{proof}

Returning departures with return index $0$ are certified. The returning departures with return index in
$\{1,\dots,k-1\}$ are paired by reversing the excursion, as in~\cite{Doz26}. Let $\mathcal T$ be the set of triples
$(\rho,i,\beta)$ with $\rho\in S$ and $(i,\beta)$ a returning departure of $\rho$ with return index
$l\in\{1,\dots,k-1\}$. For such a triple, with excursion vertices $y_0=x_i,y_1,\dots,y_r=x_l$, let $\tilde\rho$ be
the plan obtained from $\rho$ by replacing $\pi_v$ by $\pi_v^{-1}$ at every vertex $v\in\{y_1,\dots,y_{r-1}\}$, and
put
\[
\Phi(\rho,i,\beta)=(\tilde\rho,\,l,\,(x_l,y_{r-1})).
\]

\begin{lemma}\label{lem:reversal}
$\Phi$ is an involution of $\mathcal T$. For $(\rho,i,\beta)\in\mathcal T$ with $\Phi(\rho,i,\beta)=(\tilde\rho,l,\tilde\beta)$:
\begin{enumerate}
\item[(a)] $\mu(\tilde\rho)=\mu(\rho)$ and $\tilde\rho$ has the same walk $P$ from $\delta$;
\item[(b)] $Q^{-1}\le w_\rho(i)/w_{\tilde\rho}(l)\le Q$;
\item[(c)] at least one of $(i,\beta)$ for $\rho$ and $(l,\tilde\beta)$ for $\tilde\rho$ is certified.
\end{enumerate}
\end{lemma}

\begin{proof}
The derangements at the vertices of $P$ are unchanged, and the inverse of a derangement is a derangement, so
$\tilde\rho$ is a plan with walk $P$, and $\mu(\tilde\rho)=\mu(\rho)$ because $\mu$ is uniform. This is (a).

The pair $(l,\tilde\beta)$ is a departure of $\tilde\rho$. If $r\ge2$, then $y_{r-1}\notin V(P)$. If $r=1$, then
$y_{r-1}=x_i$, and $x_i=x_{l-1}$ would mean $\beta=(x_i,x_{i+1})$ while $x_i=x_{l+1}$ would mean $\beta=(x_i,x_{i-1})$,
both excluded.

For $1\le s\le r-1$ we have $\pi_{y_s}(y_{s-1})=y_{s+1}$, because $\phi_\rho(f_s)=f_{s+1}$, and the derangement of
$\tilde\rho$ at $y_s$ is $\pi_{y_s}^{-1}$, which maps $y_{s+1}$ to $y_{s-1}$. This holds also when a vertex occurs
several times among $y_1,\dots,y_{r-1}$, since the same inverse permutation is used at each occurrence. Hence
$\phi_{\tilde\rho}(y_{s+1},y_s)=(y_s,y_{s-1})$, and the excursion of $(l,\tilde\beta)$ in $\tilde\rho$ is
$(y_r,y_{r-1}),(y_{r-1},y_{r-2}),\dots,(y_1,y_0)$: its vertices $y_{r-1},\dots,y_1$ lie outside $V(P)$, and it
returns to $V(P)$ at $y_0=x_i$ after $r\le L$ edges, so its return index is $i\in\{1,\dots,k-1\}$ and
$\Phi(\rho,i,\beta)\in\mathcal T$. Applying the construction to $(\tilde\rho,l,\tilde\beta)$ inverts the same
derangements again and yields the departure $(i,(x_i,y_1))=(i,\beta)$, so $\Phi\circ\Phi$ is the identity.

For (b), $w_\rho(i)/w_{\tilde\rho}(l)=(d(x_l)-1)/(d(x_i)-1)$ lies between $(m-1)/(M-1)$ and $(M-1)/(m-1)$, and
$(M-1)/(m-1)\le Q$.

For (c), the two excursions visit the same vertices outside $V(P)$. If these visits contain a repetition, both
departures are certified. Otherwise $(i,\beta)$ is certified exactly when $l\le i$, and $(l,\tilde\beta)$ exactly when
$i\le l$.
\end{proof}

For $\rho\in S$ let $c(\rho)$ be the total weight of the certified departures of $\rho$, and let
\[
B=\frac1{\mu(S)}\sum_{\rho\in S}\mu(\rho)\,c(\rho)
\]
when $S\ne\emptyset$.

\begin{proposition}\label{prop:certified}
If $S\ne\emptyset$, then $(1+Q)B\ge\frac{k-1}2-\frac{Qn}L-Q$.
\end{proposition}

\begin{proof}
Each departure of $\rho\in S$ is certified, or non-returning and not certified, or returning with return index $k$,
or not certified with return index in $\{1,\dots,k-1\}$; returns to $x_0$ are certified. Write $u(\rho)$ for the total
weight of the last class. By~\eqref{eq:total} and Lemmas~\ref{lem:nonreturning} and~\ref{lem:terminal},
\begin{equation}\label{eq:split}
\frac{k-1}2\le c(\rho)+\frac{Qn}L+Q+u(\rho).
\end{equation}
By Lemma~\ref{lem:reversal}(c), $\Phi$ maps the triples in $\mathcal T$ whose departure is not certified injectively
into the triples in $\mathcal T$ whose departure is certified, and by (a) and (b) it changes $\mu(\rho)w_\rho(i)$ by a
factor at most $Q$. Hence
\[
\sum_{\rho\in S}\mu(\rho)u(\rho)\le Q\sum_{\rho\in S}\mu(\rho)c(\rho).
\]
Multiplying~\eqref{eq:split} by $\mu(\rho)/\mu(S)$ and summing over $\rho\in S$ gives the claim.
\end{proof}

The comparison of $u$ with $c$ holds after summation over $\rho$, not for each plan separately, since $\Phi$ changes the
plan.

\section{Transport from simple to self-intersecting plans}\label{sec:flow}

Assume $L<k$. For $\rho\in S$ let $c_0(\rho)$ be the total weight of the certified departures $(i,\beta)$ of $\rho$ with
$i\le k-L$, and let $B_0=\mu(S)^{-1}\sum_{\rho\in S}\mu(\rho)c_0(\rho)$. The departures at a single index have total
weight less than $1$, and $c_0$ omits $L-1$ indices, so $B_0\ge B-L$.

\begin{proposition}\label{prop:flow}
If $L<k$, then $p_k(\delta)\,B_0\le L$.
\end{proposition}

\begin{proof}
For $\rho\in S$, a certified departure $(i,\beta)$ of $\rho$ with $i\le k-L$ and $\beta=(x_i,z)$, and a derangement
$\sigma$ of $N(x_i)$ with $\sigma(x_{i-1})=z$, let $\psi$ be the plan obtained from $\rho$ by replacing $\pi_{x_i}$ by
$\sigma$, and send mass $\mu(\rho)/R_{d(x_i)}$ from $\rho$ to $\psi$. By Lemma~\ref{lem:certified}, $\psi\notin S$. By
Lemma~\ref{lem:derangement} there are $R_{d(x_i)}/(d(x_i)-1)$ choices of $\sigma$, so the departure $(i,\beta)$ sends
total mass $\mu(\rho)w_\rho(i)$, the plan $\rho$ sends $\mu(\rho)c_0(\rho)$, and the total mass sent is
$\mu(S)B_0=p_k(\delta)B_0$ by Lemma~\ref{lem:simple}.

Fix $\psi\notin S$ and let $j$ be the least index such that the $j$-th vertex of the walk of $\psi$ from $\delta$ equals
an earlier one. If mass is sent to $\psi$ from $(\rho,i,\beta,\sigma)$, then the walk of $\psi$ begins with the distinct
vertices $x_0,\dots,x_i$, so $j\ge i+1$, and by Lemma~\ref{lem:certified} $j\le i+L$. Thus $i$ takes at most $L$
values. Given $i$, the walk of $\psi$ determines $x_{i-1}$, $x_i$ and $\beta$, and $\sigma$ is the derangement of $\psi$
at $x_i$; the source $\rho$ agrees with $\psi$ except at $x_i$, so there are at most $R_{d(x_i)}$ sources, each sending
mass $\mu(\rho)/R_{d(x_i)}=\mu(\psi)/R_{d(x_i)}$ to $\psi$. The mass received by $\psi$ is therefore at most
$L\mu(\psi)$, and the total mass received is at most $L\mu(S^c)\le L$.
\end{proof}

Choosing the replacement derangement uniformly, rather than fixing one for each departure, is what makes the mass
received by $\psi$ from a given index at most $\mu(\psi)$; a fixed choice leads to a count of preimages that can be as
large as $R_{d(x_i)}$.

\section{Proofs of the theorem and the corollaries}\label{sec:proofs}

\begin{proof}[Proof of Theorem~\ref{thm:main}]
Put $A=1+Q$. Since $p_k(\delta)=0$ for $k\ge n$, and the bound is trivial when $k^2<256QAn$, we may assume
$1\le k\le n-1$ and $k^2\ge256QAn$, and $S\ne\emptyset$. Then $k^2\ge256QAn>256QAk$, so $k>256QA$; as $Q>1$, this gives
$k\ge16A$, $Q\le k/16$ and $k\ge4$. Let $L=\lceil16Qn/k\rceil$. From $k^2\ge256QAn$ we get $16Qn/k\le k/(16A)$, hence
\[
L\le\frac{16Qn}k+1\le\frac k{16A}+\frac k{16A}=\frac k{8A},
\]
and in particular $L<k$. Moreover $Qn/L\le k/16$. Proposition~\ref{prop:certified} gives
\[
B\ge\frac1A\Bigl(\frac{k-1}2-\frac k{16}-\frac k{16}\Bigr)\ge\frac k{4A},
\]
using $k\ge4$, so $B_0\ge B-L\ge k/(8A)$. By Proposition~\ref{prop:flow},
\[
p_k(\delta)\le\frac{L}{B_0}\le\frac{8AL}k\le\frac{128AQn}{k^2}+\frac{8A}k\le\frac{136AQn}{k^2}\le\frac{256AQn}{k^2},
\]
where the fourth inequality uses $k<n$ and $Q\ge1$.
\end{proof}

\begin{proof}[Proof of Corollary~\ref{cor:constants}]
If $M/m\le D$, then $Q=(M/m)\cdot m/(m-1)\le3D/2$, so $256Q(1+Q)\le384D+576D^2\le960D^2<1024D^2$, and
$k>32D\sqrt{n/\varepsilon}$ gives $256Q(1+Q)n/k^2<\varepsilon$. If $M\le\Delta$, then $Q\le\Delta/2$ and
$256Q(1+Q)\le64\Delta(\Delta+2)$. For regular graphs take $D=1$ in the bound $960D^2$ and use $31^2>960$.
\end{proof}

\begin{proof}[Proof of Corollary~\ref{cor:expectation}]
We have $\PP(\tau(\delta)>t)=p_t(\delta)$ for $t\ge1$, and $\PP(\tau(\delta)>0)=1$. Put $H=256Q(1+Q)n$ and
$a=\lceil\sqrt H\rceil$. By Theorem~\ref{thm:main},
\[
\E\,\tau(\delta)=\sum_{t\ge0}\PP(\tau(\delta)>t)\le a+1+H\sum_{t>a}\frac1{t^2}\le a+1+\frac Ha\le2\sqrt H+2
=32\sqrt{Q(1+Q)n}+2.
\]
The walk from $u$ whose first step is uniform on $N(u)$ is a mixture of the walks started along the edges $(u,v)$,
$v\in N(u)$, and $T(G,u)$ is the corresponding mixture of the $\tau(u,v)$; this gives the bound for $\E\,T(G,u)$. The
two special cases follow from $Q(1+Q)\le\Delta(\Delta+2)/4$ and from $Q(1+Q)\le960D^2/256<(62D/32)^2$.
\end{proof}

\begin{remark}\label{rem:ratio}
The ratio $Q$ enters the proof in three places: the packing bound of Lemma~\ref{lem:nonreturning}, which compares the
number $\sum_vd(v)$ of directed edges with the weight $1/(d(v)-1)$ of a departure; the number of departures returning to
$x_k$ in Lemma~\ref{lem:terminal}, which can be as large as $d(x_k)$ while each such departure may carry weight
$1/(m-1)$; and the change of weight under excursion reversal in Lemma~\ref{lem:reversal}(b), between a departure at a
vertex of large degree and its reversal at a vertex of small degree. The packing loss can be expressed through the
average degree instead: if a bounded number of vertices have large degree and all others have degree three, then
$\sum_vd(v)=O(n)$ and the packing bound loses only a constant factor. The terminal and reversal losses can still be
of order $M$ on such graphs, and the argument gives no bound uniform in the degrees.
Conjecture~\ref{conj:AP} without a bound on the degree ratio, and~\cite[Question~4]{BGJ26} in general, remain open.
The example after~\cite[Question~4]{BGJ26} shows that some use of simplicity of the graph is necessary, since the
analogous statements fail for multigraphs.
\end{remark}

\section*{Acknowledgements and funding}
This work was done by multiple AI agents, Claude (Anthropic) and Codex (OpenAI), connected to Math
Brain~\cite{MathBrainRepo} and working together inside the Brain framework~\cite{BrainFramework}, a shared, persistent
memory in which the agents record claims, proofs, failed attempts and methods, retrieve one another's results,
coordinate who works on what, and audit each other's work. For more on the details behind the Brain and Math Brain, see
the Math Brain paper~\cite{MathBrainPaper}, the Math Brain repository~\cite{MathBrainRepo} and the Brain framework
repository~\cite{BrainFramework}.

The research presented in this paper was supported by the European Research Council (ERC) under the European Union's
Horizon Europe research and innovation programme (grant agreement No.~101041711), by the Simons Foundation as part of
the Collaboration on the Mathematical and Scientific Foundations of Deep Learning, by Heights Labs, by Convex Nexus
Capital, by the Israel Science Foundation (grant number 2258/19), and by the Israel Science Foundation (ISF Grant
4101/25). Computational resources were provided in part by Google Cloud through its education credits programme (credit
530550103).

\section*{Declaration of competing interest}
The author holds equity in Heights Labs and in Convex Nexus Capital and has a management role at Convex Nexus Capital.
These entities are included in the funding declaration above.

\section*{Declaration of generative AI and AI-assisted technologies}
A Codex agent (OpenAI) proved Theorem~\ref{thm:main}; a second Codex agent checked the proof as a referee; Claude agents
(Anthropic) checked the steps and prepared the manuscript. The author directed the work and is responsible for the
content of the paper and for decisions concerning its dissemination.

\section*{Code availability}
A script computing $p_k(\delta)$ exactly on small random graphs and estimating the self-intersection time on larger
ones, with a comparison to Theorem~\ref{thm:main} and Corollary~\ref{cor:expectation}, is in the repository
\begin{center}\url{https://github.com/yspennstate/nonbacktracking-birthday-irregular/tree/v1.1}.\end{center}
On the graphs tested, $\max_{\delta,k}p_k(\delta)k^2/n$ is at most about $2.2$, and the mean self-intersection time is
between $1.4\sqrt n$ and $2.2\sqrt n$.

\appendix
\section{Notation}\label{app:notation}

\begin{center}
\small
\begin{tabular}{lll}
\toprule
Symbol & Meaning & First used\\
\midrule
$G=(V,E)$, $n$ & the graph and its number of vertices & Section~\ref{sec:intro}\\
$N(v)$, $d(v)$ & neighbourhood and degree of $v$ & Section~\ref{sec:intro}\\
$m$, $M$ & minimum and maximum degree & Section~\ref{sec:intro}\\
$\vE$ & directed edges & Section~\ref{sec:intro}\\
$\delta=(x_0,x_1)$ & the starting directed edge & Section~\ref{sec:intro}\\
$X_t$ & the non-backtracking walk & Section~\ref{sec:intro}\\
$p_k(\delta)$ & probability that the walk of length $k$ is vertex-simple & Section~\ref{sec:intro}\\
$T(G,u)$, $\tau(\delta)$ & self-intersection times from a vertex and from an edge & Section~\ref{sec:intro}\\
$Q$ & $M/(m-1)$ & Section~\ref{sec:intro}\\
$\varepsilon$, $C$ & a probability and the constant of Conjecture~\ref{conj:AP} & Section~\ref{sec:intro}\\
$\Delta$, $D$ & bounds for $M$ and for $M/m$ & Section~\ref{sec:intro}\\
$\PP$, $\E$ & probability and expectation for the walk & Section~\ref{sec:intro}\\
$R_d$ & number of derangements of $d$ objects & Section~\ref{sec:plans}\\
$\rho=(\pi_v)$, $\phi_\rho$ & a plan and its map on directed edges & Section~\ref{sec:plans}\\
$\mu$ & uniform measure on plans & Section~\ref{sec:plans}\\
$S$ & plans whose walk from $\delta$ is vertex-simple for $k$ steps & Section~\ref{sec:plans}\\
$P=(x_0,\dots,x_k)$, $V(P)$ & the walk of a plan in $S$ and its vertex set & Section~\ref{sec:departures}\\
$L$ & the maximal length of an excursion & Section~\ref{sec:departures}\\
$(i,\beta)$, $w_\rho(i)$ & a departure and its weight & Section~\ref{sec:departures}\\
$f_s=(y_{s-1},y_s)$, $r$, $l$ & excursion edges, length and return index & Section~\ref{sec:departures}\\
$\psi$ & a plan differing from $\rho$ at one vertex & Lemma~\ref{lem:certified}\\
$\mathcal T$, $\Phi$ & interior returns and excursion reversal & Section~\ref{sec:departures}\\
$\tilde\rho$, $\tilde\beta$ & the reversed plan and departure & Lemma~\ref{lem:reversal}\\
$c(\rho)$, $u(\rho)$, $B$ & certified weight, uncertified interior-return weight, mean of $c$ & Section~\ref{sec:departures}\\
$c_0(\rho)$, $B_0$ & the same with $i\le k-L$ & Section~\ref{sec:flow}\\
$A$, $H$ & $1+Q$ and $256Q(1+Q)n$ & Section~\ref{sec:proofs}\\
\bottomrule
\end{tabular}
\end{center}

\begin{thebibliography}{99}
\bibitem{ABLS07} N.~Alon, I.~Benjamini, E.~Lubetzky and S.~Sodin, Non-backtracking random walks mix faster,
\emph{Commun.\ Contemp.\ Math.} \textbf{9} (2007), no.~4, 585--603. \doi{10.1142/S0219199707002551}
\bibitem{BGJ26} F.~Bencs, L.~A.~Goldberg, M.~Jenssen, M.~Jerrum, G.~Pete, G.~Regts and Y.~Yin, On the self-intersection
time of non-backtracking random walks, preprint, 2026, \href{https://arxiv.org/abs/2608.09729}{arXiv:2608.09729}.
\bibitem{Doz26} B.~Dozier, The birthday paradox for non-backtracking walks on regular graphs, preprint, 2026,
\href{https://arxiv.org/abs/2608.26321}{arXiv:2608.26321}.
\bibitem{KMPT10} J.~H.~Kim, R.~Montenegro, Y.~Peres and P.~Tetali, A birthday paradox for Markov chains with an
optimal bound for collision in the Pollard rho algorithm for discrete logarithm, \emph{Ann.\ Appl.\ Probab.}
\textbf{20} (2010), no.~2, 495--521. \doi{10.1214/09-AAP625}
\bibitem{LP16} E.~Lubetzky and Y.~Peres, Cutoff on all Ramanujan graphs, \emph{Geom.\ Funct.\ Anal.}
\textbf{26} (2016), no.~4, 1190--1216. \doi{10.1007/s00039-016-0382-7}
\bibitem{MathBrainRepo} Y.~Shmalo, \emph{Math Brain: public snapshot, code and tests}, 2026,
\url{https://github.com/yspennstate/math-brain-public}.
\bibitem{MathBrainPaper} Y.~Shmalo, \emph{Math Brain: shared, auditable memory for multi-agent mathematical research},
October 2026, \url{https://github.com/yspennstate/math-brain-public/tree/main/papers/math-brain-2026-10-24}.
\bibitem{BrainFramework} Y.~Shmalo, \emph{Brain Framework}, public software repository, 2026,
\url{https://github.com/yspennstate/brain-framework-public}.
\end{thebibliography}
\end{document}
